\section*{Capítulo \ref{ch:b2:phylogenetic-trees} --- Árvores filogenéticas e relógio molecular}

\begin{solution}{exo:b2:phylogenetic-trees:1}
\omterm{def:b2:phylogenetic-trees:tree}{Homologia}: um caráter herdado de um ancestral comum --- o úmero, o rádio
e a ulna da asa de um morcego e da nadadeira de uma baleia. \omterm{def:b2:phylogenetic-trees:tree}{Analogia}: um
caráter semelhante adquirido independentemente --- as asas dos morcegos e
das aves enquanto superfícies de voo (os ossos são homólogos enquanto
membros anteriores, a asa como tal não é), ou o olho em câmara do polvo e
o dos vertebrados. \omterm{def:b2:phylogenetic-trees:tree}{Sinapomorfia}: uma \omterm{def:b2:phylogenetic-trees:tree}{homologia} derivada compartilhada por
um \omterm{def:b2:phylogenetic-trees:tree}{clado} e que o define --- as penas para as aves, o ovo amniótico para
os amniotas.
\end{solution}

\begin{solution}{exo:b2:phylogenetic-trees:2}
Aves: monofiléticas. Répteis sem as aves: parafiléticos. Peixes sem os
tetrápodes: parafiléticos (os peixes pulmonados são mais próximos de nós
que das trutas). Vertebrados voadores (morcegos, aves, pterossauros):
polifiléticos. Mamíferos: monofiléticos. Procariontes: parafiléticos
(bactérias e arqueias, com os eucariontes deixados de fora --- e as
arqueias são mais próximas dos eucariontes que das bactérias).
\end{solution}

\begin{solution}{exo:b2:phylogenetic-trees:3}
Grupo-irmão das aves: os crocodilos. Grupo-irmão dos mamíferos: todos os
outros amniotas (tartarugas, lagartos, crocodilos e aves juntos). Menor
\omterm{def:b2:phylogenetic-trees:tree}{clado} que contém lagartos e aves: (lagartos,(crocodilos, aves)), os
diápsidos sem as tartarugas nesta árvore. Trocar crocodilos e aves em seu
nó não muda nada: um nó pode ser girado livremente.
\end{solution}

\begin{solution}{exo:b2:phylogenetic-trees:4}
Não enraizadas: $3$; $3\times 5\times 7 = 105$; $3\times 5\times
7\times 9\times 11\times 13\times 15 = \num{2027025}$. Árvores
enraizadas para 4 táxons: cada uma das 3 árvores não enraizadas tem 5
ramos, o que dá 15.
\end{solution}

\begin{solution}{exo:b2:phylogenetic-trees:5}
$((A,B),(C,D))$: sítios 1, 2 e 3, um passo cada um; sítio 4, dois passos
($\mathrm{A}$ e $\mathrm{G}$ dos dois lados); sítio 5, um: 6 passos.
$((A,C),(B,D))$: $2 + 2 + 2 + 1 + 1 = 8$. $((A,D),(B,C))$: $2 + 2 + 2 +
2 + 1 = 9$. A primeira árvore é a mais parcimoniosa; nela, o sítio 4 é
homoplásico (a mudança $\mathrm{A}\to\mathrm{G}$, ou a inversa,
ocorre duas vezes).
\end{solution}

\begin{solution}{exo:b2:phylogenetic-trees:6}
Junte $A$ e $B$ na altura 3. Depois $d(AB,C) = 14$, $d(AB,D) = 16$,
$d(C,D) = 10$: junte $C$ e $D$ na altura 5. Depois $d(AB,CD) = (14 + 16
+ 14 + 16)/4 = 15$: raiz na altura $7.5$. Árvore $((A,B),(C,D))$.
\end{solution}

\begin{solution}{exo:b2:phylogenetic-trees:7}
$d = -0.75\ln(1 - 4p/3)$: $0.052$, $0.38$, $1.21$. Com $p = 0.6$, o
argumento do logaritmo é $0.2$ e a inclinação $\mathrm{d}d/\mathrm{d}p
= 1/(1 - 4p/3) = 5$: um erro de amostragem de $0.02$ em $p$ se torna
$0.1$ em $d$, e a hipótese de taxas iguais em todos os sítios, falsa em
todo gene real, acrescenta um viés do mesmo tamanho. O gene está
saturado.
\end{solution}

\begin{solution}{exo:b2:phylogenetic-trees:8}
$t = d/2k = 0.02/(2\times 10^{-9}) = 10$ milhões de anos. A 500
milhões de anos, $d = 2\times 10^{-9}\times 5\times 10^{8} = 1.0$ e
$p = 0.75(1 - e^{-4/3}) = 0.55$: três quartos do caminho até a
saturação, onde a correção é íngreme e a data não é confiável --- é
preciso um gene mais lento para essas separações.
\end{solution}

\begin{solution}{exo:b2:phylogenetic-trees:9}
\qty{52}{\%}: o \omterm{def:b2:phylogenetic-trees:tree}{clado} é recuperado em apenas metade das reamostragens;
os dados quase nada dizem sobre ele, e as alternativas, somadas, são
igualmente prováveis. \qty{99}{\%}: o sinal é coerente entre os sítios;
o \omterm{def:b2:phylogenetic-trees:tree}{clado} é fortemente sustentado (dados o modelo e o alinhamento --- os
erros sistemáticos, como a \omterm{prop:b2:phylogenetic-trees:pitfalls}{atração de ramos longos}, não são detectados
pelo bootstrap). Para o primeiro, acrescente sequência (mais genes),
acrescente táxons que quebrem os ramos e verifique com um método
diferente.
\end{solution}

\begin{solution}{exo:b2:phylogenetic-trees:10}
Duas linhagens que mudaram, cada uma, numa grande fração dos sítios têm,
em cada sítio onde ambas mudaram, uma chance em quatro de chegar à mesma
base (com quatro bases a taxas iguais). Se cada uma mudou em
\qty{40}{\%} dos sítios, ambas mudaram em \qty{16}{\%} e coincidem por
acaso em \qty{4}{\%} de todos os sítios --- o que pode superar a fração
de sítios que carregam as verdadeiras \omterm{def:b2:phylogenetic-trees:tree}{sinapomorfias} que ligam cada uma a
seus parentes reais, de evolução lenta. A parcimônia conta coincidências
sem perguntar por que elas ocorrem e junta as duas. Acrescentar táxons
que se ramificam ao longo de cada ramo longo o divide em segmentos mais
curtos, de modo que as coincidências fortuitas se distribuem entre
vários ramos mais curtos e as verdadeiras \omterm{def:b2:phylogenetic-trees:tree}{sinapomorfias} de cada segmento
se tornam visíveis; um modelo de verossimilhança que espera muitas
coincidências fortuitas nos ramos longos as corrige diretamente.
\end{solution}

\begin{solution}{exo:b2:phylogenetic-trees:11}
$d = -0.75\ln(1 - 0.12) = 0.096$; $t = d/2k = 0.096/(4\times 10^{-8})
= 2.4$ milhões de anos. Substituições: $0.096\times \num{16000} =
1500$, precisão relativa $1/\sqrt{1500} = \qty{2.6}{\%}$: um erro
estatístico de $\pm 0.06$ milhão de anos. Mas a taxa é calibrada em
outras separações e pode estar errada por um fator de dois para esta
linhagem (variação das taxas entre linhagens, e diferença entre a taxa
genealógica medida ao longo de gerações e a taxa medida ao longo de
milhões de anos); e a divergência do gene é anterior à divergência das
espécies, pois a população ancestral era polimórfica --- de modo que a
verdadeira separação é mais recente que a do gene, por um tempo que
depende do tamanho da população ancestral. (A data aceita, a partir de
muitos genes nucleares e de fósseis, é de 6 a 7 milhões de anos: o
relógio mitocondrial aqui é rápido demais.)
\end{solution}

\begin{solution}{exo:b2:phylogenetic-trees:12}
\omterm{prop:b2:phylogenetic-trees:pitfalls}{Separação incompleta de linhagens}: quando duas especiações são próximas
no tempo, um polimorfismo ancestral pode se distribuir de modo que a
árvore de um gene agrupe as espécies de forma diferente da árvore de
espécies; transferência horizontal: um gene recebido de outra linhagem
carrega a história dessa linhagem; \omterm{prop:b2:genome-diversification:duplication}{duplicação gênica}: comparar um
parálogo numa espécie com o ortólogo em outra dá a data da duplicação, e
não a da especiação. Obtém-se uma árvore de espécies usando muitos genes
de todo o genoma e tomando a árvore que a maioria sustenta (ou um modelo
que espera uma fração conhecida de genes discordantes), escolhendo
cuidadosamente os ortólogos e preferindo, nos procariontes, os genes
ribossômicos e informacionais, que raramente são transferidos.
\end{solution}

\begin{solution}{pb:b2:phylogenetic-trees:1}
\textbf{1.} Sítios 1, 2 e 4 (estados compartilhados apenas por $W$ e $H$,
derivados em relação ao camelo): \omterm{def:b2:phylogenetic-trees:tree}{sinapomorfias} de $\{W,H\}$. Sítios 3
e 6: \omterm{def:b2:phylogenetic-trees:tree}{sinapomorfias} de $\{W,H,C\}$. O sítio 5 é uma autapomorfia de
$W$, uma mudança numa única linhagem, que não agrupa nada; nenhum sítio é
homoplásico na árvore molecular.
\textbf{2.} Uma mudança por sítio: 6 passos.
\textbf{3.} Os sítios 1, 2 e 4 precisam agora de duas mudanças cada um,
e os sítios 3, 5 e 6 de uma: 9 passos.
\textbf{4.} A árvore baleia--hipopótamo, por três passos. Três sítios em
mil é uma margem pequena, que algumas homoplasias poderiam inverter;
reforçá-la significa mais genes, mais táxons (outros ruminantes,
caititus, cervos) e um bootstrap.
\textbf{5.} O \omterm{def:b2:phylogenetic-trees:tree}{clado} reaparece em cinco reamostragens de cada seis: um
suporte moderado, bem acima do acaso, mas aquém dos \qty{95}{\%}
convencionalmente chamados de forte.
\textbf{6.} O \omterm{meth:b2:phylogenetic-trees:parsimony}{grupo externo} deve estar fora do grupo, mas perto o
bastante para que sua sequência ainda seja alinhável e não saturada; o
camelo é um artiodáctilo situado fora do grupo vaca--porco--hipopótamo--baleia. Um
peixe diferiria em níveis quase saturados, seus estados não
informariam qual estado é ancestral dentro dos mamíferos,
e seu ramo longo atrairia a linhagem do grupo interno de evolução mais
rápida.
\textbf{7.} $d = 0.0625$, $0.107$, $0.131$, $0.155$.
\textbf{8.} Junte $W$ e $H$ na altura $0.031$. Depois $d(WH,C) =
0.107$, $d(WH,P) = 0.131$, $d(C,P) = 0.131$, e tudo a $M$
$0.155$: junte $WH$ e $C$ em $0.054$. Depois $d(WHC,P) = 0.131$: junte
em $0.065$. Por fim, junte $M$ em $0.078$. Árvore $(M,(P,(C,(W,H))))$.
\textbf{9.} Sim: a mesma árvore, o mesmo encaixe.
\textbf{10.} Uma baleia que tivesse mudado duas vezes mais depressa
estaria mais longe de todos, inclusive do hipopótamo; o UPGMA, que coloca
cada ponta na altura zero e junta o par mais próximo, juntaria primeiro o
hipopótamo com a vaca e deixaria a baleia de fora --- exatamente a
árvore dos anatomistas, pela razão errada.
\textbf{11.} $0.0625\times 1000 = 62.5$ substituições; desvio-padrão
de Poisson $\sqrt{62.5} = 7.9$.
\textbf{12.} $7.9/62.5 = \qty{13}{\%}$.
\textbf{13.} $k = d/2t = 0.107/(1.2\times 10^{8}) = 8.9\times 10^{-10}$
por sítio por ano.
\textbf{14.} $t = 0.0625/(2\times 8.9\times 10^{-10}) = 35$ milhões
de anos.
\textbf{15.} Porco: $0.131/(1.79\times 10^{-9}) = 73$ milhões de anos;
camelo: $0.155/(1.79\times 10^{-9}) = 87$ milhões de anos.
\textbf{16.} $35 \pm 4.5$ milhões de anos (\qty{13}{\%}).
\textbf{17.} A taxa varia como $1/t_{\text{cal}}$ e a data como
$t_{\text{cal}}$: $\pm 5/60 = \qty{8}{\%}$, logo $\pm 3$ milhões de anos
--- combinado com o erro de Poisson, cerca de $\pm 5$ milhões de anos.
\textbf{18.} $d = -0.75\ln(1 - 0.6) = 0.69$: o gene mudou, na prática,
em dois terços de seus sítios, a correção multiplicou
$p$ por $1.5$ e sua inclinação é $2.5$; ele está perto da saturação para
essa separação, e sua data vale pouco.
\textbf{19.} A ausência de um caráter não é um estado derivado
compartilhado: as baleias perderam inteiramente o membro posterior e seu
tornozelo, de modo que nada se pode dizer sobre qual tornozelo teriam
tido. Uma \omterm{def:b2:phylogenetic-trees:tree}{sinapomorfia} é evidência de um \omterm{def:b2:phylogenetic-trees:tree}{clado}; sua perda numa linhagem
não é evidência contra o pertencimento.
\textbf{20.} As baleias fósseis têm o astrágalo em dupla polia: as
baleias descendem de um ancestral que o possuía, o que as coloca dentro
dos artiodáctilos. A previsão molecular é confirmada pelas
rochas.
\textbf{21.} Parafilético: um ancestral com todos os seus descendentes,
exceto a linhagem das baleias.
\textbf{22.} \omterm{prop:b2:phylogenetic-trees:pitfalls}{Separação incompleta de linhagens} de um polimorfismo
ancestral ao longo de duas especiações rápidas (vaca, hipopótamo e baleia
se separaram em poucos milhões de anos), ou comparação de parálogos.
Decide-se sequenciando muitos genes e tomando a árvore que a maioria
deles, e a concatenação, sustentam; verifica-se a duplicação em cada
família.
\textbf{23.} $d = -0.75\ln(1 - 0.333) = 0.30$; $k = d/2t =
0.30/(7\times 10^{7}) = 4.3\times 10^{-9}$ por sítio por ano, cinco
vezes a do gene nuclear.
\textbf{24.} O DNA mitocondrial é replicado por uma polimerase de
revisão mais fraca, exposto aos radicais de oxigênio da cadeia
respiratória e não é reparado pela maquinaria nuclear; ele também não
tem recombinação, de modo que os danos se acumulam.
\textbf{25.} Árvore $(M,(P,(C,(W,H))))$: as baleias são o grupo-irmão
dos hipopótamos; $k = 8.9\times 10^{-10}$ por sítio por ano; divergência
baleia--hipopótamo de $35 \pm 5$ milhões de anos.
\end{solution}
